% ============================================================
% 学术汇报 Beamer 模板
% 编译: xelatex beamer-slides.tex
% ============================================================
\documentclass[aspectratio=169,11pt]{beamer}

% === 主题 ===
\usetheme{Madrid}
\usecolortheme{default}
\setbeamertemplate{navigation symbols}{}
\setbeamertemplate{footline}[frame number]
\setbeamertemplate{caption}[numbered]

% === 包 ===
\usepackage{amsmath,amssymb}
\usepackage{graphicx}
\usepackage{booktabs}
\usepackage{tikz}
\usepackage{appendixnumberbeamer}

% === 中文支持（如需要） ===
% \usepackage{xeCJK}
% \setCJKmainfont{SimSun}

% === 自定义颜色 ===
\definecolor{mainblue}{RGB}{44,62,80}
\definecolor{accentred}{RGB}{192,57,43}
\setbeamercolor{title}{fg=mainblue}
\setbeamercolor{frametitle}{fg=mainblue}
\setbeamercolor{structure}{fg=mainblue}

% === 标题信息 ===
\title[\shorttitle]{[论文完整标题]}
\newcommand{\shorttitle}{[短标题]}
\author{[作者]}
\institute{[单位]}
\date{[日期/会议名称]}

\begin{document}

% ============================================================
\begin{frame}
\titlepage
\end{frame}
% ============================================================

% ============================================================
\begin{frame}{Outline}
\tableofcontents
\end{frame}
% ============================================================


% ============================================================
\section{Motivation}
% ============================================================

\begin{frame}{Research Question}
\begin{itemize}
    \item \textbf{Core Question:} [一句话研究问题]
    \vspace{1em}
    \item \textbf{Why it matters:}
    \begin{itemize}
        \item [原因1 — 政策/实践意义]
        \item [原因2 — 学术空白]
    \end{itemize}
    \vspace{1em}
    \item \textbf{What we find:} [一句话核心发现]
\end{itemize}
\end{frame}

\begin{frame}{Contribution}
\begin{enumerate}
    \item \textbf{[贡献1]:} [简要说明]
    \item \textbf{[贡献2]:} [简要说明]
    \item \textbf{[贡献3]:} [简要说明]
\end{enumerate}

\vspace{1em}
\textbf{Related literature:} [Author1 (Year)]; [Author2 (Year)]; [Author3 (Year)]
\end{frame}


% ============================================================
\section{Background}
% ============================================================

\begin{frame}{Institutional Background}
\begin{columns}
\column{0.55\textwidth}
\begin{itemize}
    \item [政策/制度描述]
    \item [关键时间节点]
    \item [适用范围]
\end{itemize}

\column{0.4\textwidth}
% 可放政策时间线或示意图
\begin{center}
\textbf{Timeline}\\[0.5em]
\begin{tabular}{ll}
\toprule
Date & Event \\
\midrule
[日期] & [事件1] \\
[日期] & [事件2] \\
[日期] & [事件3] \\
\bottomrule
\end{tabular}
\end{center}
\end{columns}
\end{frame}


% ============================================================
\section{Data \& Strategy}
% ============================================================

\begin{frame}{Data}
\begin{itemize}
    \item \textbf{Source:} [数据来源]
    \item \textbf{Sample:} [样本描述]
    \item \textbf{Period:} [时间范围]
    \item \textbf{Key variables:}
    \begin{itemize}
        \item $Y$: [因变量]
        \item $D$: [处理变量]
        \item Controls: [控制变量]
    \end{itemize}
\end{itemize}
\end{frame}

\begin{frame}{Empirical Strategy}
\textbf{[方法名称]} (e.g., Difference-in-Differences)

\vspace{0.5em}
\begin{equation*}
Y_{it} = \alpha_i + \delta_t + \beta(\text{Treat}_i \times \text{Post}_t) + X'_{it}\gamma + \varepsilon_{it}
\end{equation*}

\begin{itemize}
    \item $\alpha_i$: entity fixed effects
    \item $\delta_t$: time fixed effects
    \item $\beta$: \textcolor{accentred}{\textbf{coefficient of interest}} — [经济含义]
    \item SE clustered at [层级] level
\end{itemize}

\vspace{0.5em}
\textbf{Identification assumption:} [平行趋势/排他性等]
\end{frame}


% ============================================================
\section{Results}
% ============================================================

\begin{frame}{Main Results}
\begin{table}
\centering
\small
\caption{[表标题]}
\begin{tabular}{lccc}
\toprule
 & (1) & (2) & (3) \\
\midrule
Treat $\times$ Post & [coef]*** & [coef]*** & [coef]*** \\
 & ([se]) & ([se]) & ([se]) \\
\midrule
Controls & No & Yes & Yes \\
Entity FE & Yes & Yes & Yes \\
Year FE & Yes & Yes & Yes \\
Obs. & [N] & [N] & [N] \\
\bottomrule
\end{tabular}
\end{table}
{\small \textit{Notes: Clustered SE in parentheses. *** $p<0.01$}}
\end{frame}

\begin{frame}{Event Study / Parallel Trends}
\begin{figure}
\centering
\includegraphics[width=0.8\textwidth]{figures/fig2_event_study.png}
\caption{Dynamic Treatment Effects}
\end{figure}
{\small Pre-trend coefficients jointly insignificant ($F = [X.XX]$, $p = [X.XX]$)}
\end{frame}

\begin{frame}{Robustness}
\begin{itemize}
    \item \textcolor{green!50!black}{\checkmark} \textbf{Placebo test:} [结果]
    \item \textcolor{green!50!black}{\checkmark} \textbf{Alternative control group:} [结果]
    \item \textcolor{green!50!black}{\checkmark} \textbf{Alternative DV:} [结果]
    \item \textcolor{green!50!black}{\checkmark} \textbf{Different clustering:} [结果]
\end{itemize}
\vspace{0.5em}
$\Rightarrow$ Results robust across all specifications
\end{frame}

\begin{frame}{Heterogeneity}
\begin{columns}
\column{0.5\textwidth}
\textbf{By [维度1]:}
\begin{itemize}
    \item [子组A]: $\beta = [X]$***
    \item [子组B]: $\beta = [X]$ (n.s.)
\end{itemize}

\column{0.5\textwidth}
\textbf{By [维度2]:}
\begin{itemize}
    \item [子组C]: $\beta = [X]$***
    \item [子组D]: $\beta = [X]$*
\end{itemize}
\end{columns}

\vspace{1em}
\textbf{Interpretation:} [异质性的经济解释]
\end{frame}


% ============================================================
\section{Conclusion}
% ============================================================

\begin{frame}{Conclusion}
\begin{enumerate}
    \item \textbf{Finding:} [核心发现]
    \vspace{0.5em}
    \item \textbf{Mechanism:} [机制/渠道]
    \vspace{0.5em}
    \item \textbf{Policy implication:} [政策建议]
    \vspace{0.5em}
    \item \textbf{Limitation:} [主要局限]
\end{enumerate}

\vspace{1em}
\centering
\Large Thank you! \\
\normalsize [邮箱]
\end{frame}


% ============================================================
% 附录
% ============================================================
\appendix

\begin{frame}[noframenumbering]{Appendix: Additional Results}
[备用幻灯片，应对提问]
\end{frame}

\end{document}
